\documentclass[10pt,a4paper]{article}
\usepackage[margin=20mm]{geometry}
\usepackage{amsmath,amssymb,hyperref,microtype}
\hypersetup{colorlinks=true,linkcolor=blue,urlcolor=blue}
\title{Report 29 - ElasticPenstock Model}
\author{Madhusudhan Pandey}
\date{20 September 2026}
\begin{document}\maketitle
\begin{abstract}This report documents the ElasticPenstock model family in OpenHPLjl. Catalogue status: \textbf{prototype}.\end{abstract}
\section{Updated Engineering Context}Compressible elastic penstock family. It belongs to the Rigid pipes and waterways domain.
\section{Generalized Concept Model}The model preserves the domain connector interface while allowing internal fidelity to evolve independently.
\section{Alternative Models Available in Literature}This family sits on the common fidelity ladder: boundary/algebraic, lumped dynamic, nonlinear physics-based, and distributed/high-order.
\section{Governing Mathematics}\[\partial_t H+\frac{a^2}{gA}\partial_x Q=0\]The production model may add states, constitutive relations, lookup maps, or spatial discretization.
\section{OpenHPL / Modelica Interpretation}Report 9 maps this family to the closest OpenHPL, Modelica Standard Library, or OpenIPSL analogue.
\section{Julia / ModelingToolkit Implementation}Source: \texttt{\detokenize{src/Waterways/RigidPipes/ElasticPenstock.jl}}. The file is either an implemented component or an explicit Julia source specification.
\section{Component Connections and Assembly}The model connects through the stable physical connector for its domain and should not expose internal fidelity choices to neighbouring components.
\section{Numerical Experiment / Validation}Promotion requires an equation-level test, component experiment, limiting-case comparison, structural check, and connected-system test.
\section{Expected Outputs}Expected outputs are the physical states, connector variables, and conservation or residual quantities needed to verify this family.
\section{Engineering Interpretation}Choose a semi-discretization and validate wave travel time.
\section*{Conclusion}ElasticPenstock now has an explicit source/report pair. Promotion to baseline requires validated equations, parameters, tests, and reference behaviour.
\end{document}